\chapter{Fundamentação Teórica}
\label{cap:fundamentacao}

Este capítulo reúne os conceitos que sustentam as decisões metodológicas do
trabalho: o funcionamento do \textit{Bounded Model Checking}, as razões para o
uso do ESBMC em firmware embarcado e a arquitetura do PX4 Autopilot, alvo da
verificação.

\section{Verificação Formal de Software}

A verificação formal emprega métodos matemáticos para avaliar se um modelo do
sistema satisfaz propriedades previamente especificadas~\cite{Clarke1986}.
Diferentemente dos testes convencionais, que examinam o comportamento em cenários
selecionados, essas técnicas analisam sistematicamente os comportamentos
representados pelo modelo, de acordo com as hipóteses e os limites estabelecidos.

Este trabalho concentra-se em propriedades de \textit{safety} --- condições que
não devem ocorrer, como acessos fora dos limites de \textit{arrays} ou
desreferência de ponteiros nulos. A escolha decorre do domínio: em firmware
embarcado sem proteção de memória robusta, uma violação de \textit{safety} pode
corromper o estado global e causar falha catastrófica em voo. Além disso, o
ESBMC verifica automaticamente diversas propriedades de \textit{safety} --- como
limites de \textit{arrays} e validade de ponteiros --- e permite que a
pesquisadora especifique propriedades adicionais por meio de instruções
\texttt{assert}, de modo que a análise cubra tanto condições de segurança de
memória quanto comportamentos próprios de cada módulo. Propriedades de
\textit{liveness} --- que garantem que algo desejável acabe ocorrendo --- ficaram
fora do escopo: o suporte do ESBMC a elas limita-se à terminação e à detecção de
laços não terminantes, sem verificação geral de LTL~\cite{Menezes2024}.
%% LC (resolvido): afirmacao sobre liveness recalibrada ao escopo real do ESBMC
%% (terminacao/lacos nao terminantes, nao LTL geral); citacao Clarke1986 -> Menezes2024.

%% LC (resolvido): DART (Godefroid2005) e teste concolico; execucao simbolica passou
%% a citar o artigo original, King1976.
Entre as técnicas de verificação formal destacam-se o \textit{Model
Checking}~\cite{Clarke1981}, que explora automaticamente o espaço de estados; o
\textit{Theorem Proving}~\cite{Huth2004,Tang2025}, que constrói provas com auxílio de
assistentes de prova; a \textit{Abstract Interpretation}~\cite{Cousot1977}, que
aproxima o comportamento por domínios abstratos; e a \textit{Symbolic
Execution}~\cite{King1976}, que executa o programa com valores simbólicos. Pela
automação e pela adequação a software embarcado, adota-se aqui o Model Checking.

\section{Model Checking}

O Model Checking, introduzido por Clarke e Emerson~\cite{Clarke1981} e, de forma
independente, por Queille e Sifakis~\cite{Queille1982}, constrói um modelo
matemático do sistema --- tipicamente uma máquina de estados finita --- e
verifica algoritmicamente se ele satisfaz uma especificação em lógica
temporal~\cite{Clarke1986}. Como esse modelo é uma abstração do sistema real, a
validade da conclusão depende da fidelidade da representação construída.

\subsection{Processo de Model Checking}

O processo tem três etapas. Na modelagem, o sistema é representado como uma
estrutura de Kripke, em que os vértices são estados e as arestas, transições. As
propriedades são escritas em lógicas temporais como LTL (\textit{Linear Temporal
Logic})~\cite{Pnueli1977} ou CTL (\textit{Computation Tree Logic})~\cite{Clarke1986};
uma propriedade LTL típica, $\square(\textit{request} \rightarrow
\Diamond\textit{response})$, lê-se ``toda requisição é eventualmente
respondida'', com $\square$ denotando ``sempre'' e $\Diamond$, ``eventualmente''.
Por fim, o verificador percorre os estados e, se algum violar a propriedade,
retorna um contraexemplo --- a sequência de transições que leva à
violação~\cite{Clarke1986}. Essa geração automática de contraexemplos é o que
torna o Model Checking atraente para sistemas complexos.

%% LC (resolvido): algoritmo restrito a invariantes/safety -- renomeado e com
%% ressalva de que propriedades temporais gerais exigem rotulacao CTL ou automato
%% de Buchi (Clarke1986). LC (resolvido): acrescentado o mapa de predecessores.
O Algoritmo~\ref{alg:mc} apresenta, de forma simplificada, a busca de invariantes
no espaço de estados; é uma ilustração do princípio, e não o procedimento
completo das demais verificações temporais. Propriedades temporais gerais não se
decidem olhando um estado isolado: exigem a rotulação de CTL ou a busca de um
ciclo aceitante em um autômato de Büchi~\cite{Clarke1986}. O caso aqui ilustrado
--- \textit{safety}/invariantes --- é o empregado nesta pesquisa.

\begin{algorithm}
\caption{Algoritmo de verificação de invariantes por busca no espaço de estados}
\label{alg:mc}
\begin{algorithmic}[1]
\Require Modelo $M = (S, S_0, T, L)$, propriedade $\phi$
\Ensure \textit{verdadeiro} se $M \models \phi$, contraexemplo caso contrário
\State $estados \leftarrow S_0$
\State $visitados \leftarrow \emptyset$
\While{$estados \neq \emptyset$}
    \State $s \leftarrow \text{remover}(estados)$
    \If{$s \notin visitados$}
        \State $visitados \leftarrow visitados \cup \{s\}$
        \If{$s \not\models \phi$}
            \State \Return contraexemplo até $s$
        \EndIf
        \State $estados \leftarrow estados \cup \text{sucessores}(s, T)$
    \EndIf
\EndWhile
\State \Return \textit{verdadeiro}
\end{algorithmic}
\end{algorithm}

O conjunto $estados$ representa os estados ainda não explorados; a estrutura usada
para selecionar o próximo elemento determina a estratégia de busca --- uma fila
leva à busca em largura, uma pilha à busca em profundidade ---, e ambas preservam
a corretude para invariantes, desde que todos os estados alcançáveis sejam
explorados. Na prática, mantém-se um mapa de predecessores durante a busca: ao
encontrar um estado violador, o contraexemplo é reconstruído seguindo esse mapa,
do estado inicial até ele.

\subsection{Explosão de Estados}

O principal obstáculo do Model Checking é a explosão combinatória: o número de
estados cresce exponencialmente com o número de variáveis ($2^n$ para $n$
variáveis booleanas). O driver do IMU BMI088, por exemplo, manipula leituras
brutas do sensor e valores convertidos em ponto flutuante de 32 bits; um modelo que
tratasse como simbólicas apenas dez dessas variáveis de 32 bits já ultrapassaria
$2^{320}$ configurações, inviabilizando a exploração exaustiva.
%% LC (resolvido): hipotese do 2^320 declarada (dez variaveis de 32 bits).
Diversas técnicas atacam esse problema ---
BDDs (\textit{Binary Decision Diagrams}), abstração, redução de ordem parcial e o
\textit{Bounded Model Checking}, que limita a profundidade de
exploração~\cite{Biere1999}.

\section{Bounded Model Checking}

O \textit{Bounded Model Checking} (BMC), proposto por Biere et
al.~\cite{Biere1999}, restringe a busca por contraexemplos a execuções de até $k$
transições, com $k$ definido pelo usuário, em vez de explorar todos os estados
alcançáveis.

\subsection{Formulação de BMC}

Dado o estado inicial $I$, a relação de transição $T$ e uma propriedade de
\textit{safety} $\phi$, o BMC codifica a busca por contraexemplos como um
problema de satisfatibilidade:

\begin{equation}
BMC_{\phi}(k) = I(s_0) \wedge \bigwedge_{i=0}^{k-1} T(s_i, s_{i+1}) \wedge
\bigvee_{i=0}^{k} \neg\phi(s_i)
\label{eq:bmc}
\end{equation}

O termo $I(s_0)$ fixa as condições iniciais; a conjunção central codifica $k$
transições válidas; e a disjunção final verifica se algum estado $s_i$ viola
$\phi$. Se a fórmula for satisfatível, o solver devolve uma atribuição que
constitui um contraexemplo concreto~\cite{Biere1999}.

%% LC (resolvido): incluir exemplo completo (laco C -> desenrolamento -> SSA ->
%% VCC -> consulta SMT). LC (resolvido): relacionar a formula do BMC a codificacao
%% do ESBMC para programas C -- ambos atendidos pelo exemplo abaixo.
Em um programa C, essa codificação parte do código real. Considere um laço que
percorre um \textit{array} de $N$ posições (Listagem~\ref{lst:exemplo-bmc}):

\begin{lstlisting}[caption={Laço em C cuja segurança de acesso o ESBMC verifica.},
  label={lst:exemplo-bmc}, numbers=left, frame=single, language=C]
int soma = 0;
for (int i = 0; i < n; i++)
    soma += v[i];      // acesso a v[i]: vale 0 <= i < N ?
\end{lstlisting}

O vetor \texttt{v} tem \texttt{N} posições, e o laço acessa \texttt{v[i]} para
\texttt{i} de $0$ a \texttt{n}$-1$. Considerando \texttt{n} não negativo, a
condição de segurança é que o número de iterações não ultrapasse o tamanho do
vetor, isto é, $0 \leq n \leq N$; para que a análise seja conclusiva, o limite de
desenrolamento também deve ser suficiente para representar todas as iterações
possíveis. O ESBMC desenrola o laço $k$ vezes, substituindo cada iteração por uma
cópia condicionada a \texttt{i < n}; converte o resultado à forma SSA
(\textit{Static Single Assignment}), em que cada atribuição gera uma nova versão da
variável ($soma_1, soma_2, \dots$); e, para cada acesso \texttt{v[i]}, gera uma
condição de verificação (VCC) de limite ($0 \leq i < N$). As VCCs são submetidas ao
solver SMT: se \texttt{n} puder exceder \texttt{N}, a negação da condição de
segurança pode ser satisfatível, permitindo ao solver produzir uma atribuição que
representa o acesso fora dos limites do vetor; caso contrário, o resultado é UNSAT
até $k$.

\subsection{Tratamento de Laços, Completude e $k$-Indução}

Aplicar BMC a programas exige desenrolar (\textit{unwinding}) os laços até $k$.
Uma \textit{assertion} ao fim de cada laço verifica se ele poderia executar além
de $k$; quando ela falha, a análise daquele laço deixa de ser conclusiva, e é
preciso elevar $k$ ou recorrer à $k$-indução --- foi o que ocorreu com o MAVLink,
verificado com $k = 300$ (Capítulo~\ref{cap:experimentos}). Com a opção
\texttt{--no-unwinding-assertions}, laços acima de $k$ são truncados
silenciosamente~\cite{Gadelha2017}.
%% LC (resolvido): Biere1999 nao sustenta o comportamento da flag do ESBMC;
%% citacao trocada para Gadelha2017.

Sob uma codificação fiel do modelo e com o limite de desenrolamento adequadamente
tratado, o BMC é \textit{sound} e completo até $k$: todo contraexemplo encontrado
representa uma violação da propriedade \emph{no modelo}, e toda violação alcançável
em até $k$ passos é identificada; a ausência de contraexemplos, porém, indica
apenas que nenhuma violação
foi encontrada dentro do limite de exploração, e não demonstra, isoladamente, que
a propriedade se mantém em execuções que ultrapassem esse limite~\cite{Biere1999}. A correspondência entre esse contraexemplo e o código de
produção não é automática --- depende da fidelidade do \textit{harness} ---, e é o
que a etapa de conferência (etapa~(v), Capítulo~\ref{cap:metodologia})
estabelece. Para estender a garantia além de $k$, a $k$-indução~\cite{Sheeran2000}
combina um caso base (a propriedade vale nos primeiros $k$ estados) com um passo
indutivo (se vale em $k$ estados consecutivos quaisquer, vale no seguinte); o
ESBMC a implementa automaticamente~\cite{Gadelha2017}. A propriedade só é
considerada comprovada quando o caso base e o passo indutivo são ambos verificados
com sucesso. Distinta dessa $k$-indução automática, a prova de inalcançabilidade do GPS uBlox apresentada nesta
dissertação (Seção~\ref{subsec:gps-inducao}) é uma indução manual de um único
passo sobre um invariante do laço.
%% LC (resolvido): k-inducao e de Sheeran2000 (nao Biere1999), implementada no
%% ESBMC por Gadelha2017; distincao entre inducao manual de um passo e a k-inducao
%% automatica da ferramenta acrescentada.

\section{Satisfiability Modulo Theories}
\label{sec:smt}

O \textit{Satisfiability Modulo Theories} (SMT) generaliza a satisfatibilidade
booleana (SAT) ao raciocinar sobre teorias matemáticas
específicas~\cite{Barrett2018}. Os solvers SMT combinam procedimentos de decisão
para teorias como aritmética, \textit{arrays} e \textit{bit-vectors} com
algoritmos SAT eficientes, resolvendo fórmulas impraticáveis em lógica
proposicional pura. Entre os solvers atuais estão o Z3~\cite{deMoura2008}, o
Bitwuzla~\cite{Niemetz2023} e o cvc5~\cite{Barbosa2022}; o
Boolector~\cite{Brummayer2009}, especializado em \textit{bit-vectors} e
\textit{arrays}, é o solver usado em todas as verificações deste trabalho
(versão~3.2.4).
%% LC (resolvido): definicao de SMT estava sem citacao -- acrescentada a
%% referencia padrao Barrett2018 (Handbook of Model Checking, 2018).

\subsection{Teorias Relevantes e Vantagens sobre SAT}

A teoria de \textit{bit-vectors} modela inteiros de largura fixa, capturando a
semântica de tipos como \texttt{int32\_t} e \texttt{uint8\_t}, inclusive
\textit{overflow} e \textit{underflow}. A de \textit{arrays} representa estruturas
indexadas por operações de leitura (\textit{select}) e escrita (\textit{store}),
essencial para modelar memória e \textit{buffers}. Para dados de sensores, a
teoria de ponto flutuante implementa a aritmética IEEE-754, com valores especiais
como NaN e infinito. O Boolector, porém, não possui teoria nativa de ponto
flutuante: o ESBMC contorna isso codificando cada operação IEEE-754 em operações
de \textit{bit-vector} (\textit{bit-blasting}), de modo que solvers restritos a
\textit{bit-vectors}, como o Boolector, verifiquem programas com ponto
flutuante~\cite{Gadelha2020}.
%% LC (resolvido): Boolector nao tem teoria nativa de ponto flutuante --
%% explicada a codificacao por bit-blasting do ESBMC (Gadelha2020).

Frente à codificação direta em SAT, os solvers SMT são mais expressivos ---
representam restrições como $x + y > 10$ sem codificação bit a bit --- e exploram
a estrutura matemática das teorias, ganhando eficiência~\cite{deMoura2008}. No
PX4, essa expressividade é essencial: os módulos operam com ponto flutuante
IEEE-754, inteiros sem sinal sujeitos a \textit{underflow} e \textit{arrays} de
tamanho fixo para \textit{buffers} FIFO --- classes que os solvers SMT tratam com
precisão (o ponto flutuante, no caso do Boolector, pela codificação em
\textit{bit-vectors} descrita acima).

\section{ESBMC}

O ESBMC (\textit{Efficient SMT-Based Context-Bounded Model Checker}) verifica
programas em C e C++, com \textit{front-ends} adicionais para
Python~\cite{Farias2024}, Kotlin~\cite{Menezes2022} e
Solidity~\cite{Menezes2024}. Mantido por uma colaboração internacional que inclui
pesquisadores da Universidade Federal do Amazonas, adota uma abordagem baseada em
SMT, combinando a eficiência dos procedimentos SAT com o raciocínio sobre teorias
matemáticas --- aritmética, \textit{bit-vectors} e \textit{arrays} ---, o que dá
suporte mais preciso à verificação de aritmética e estruturas de dados
complexas~\cite{Cordeiro2012}.
%% LC (resolvido): sigla ESBMC uniformizada para "Efficient SMT-Based
%% Context-Bounded Model Checker" nos quatro lugares do documento.

\subsection{Pipeline de Verificação}

O fluxo do ESBMC tem quatro estágios (Figura~\ref{fig:esbmc_pipeline}). O
código-fonte passa pelo \textit{front-end} da linguagem --- Clang para C/C++,
Soot para Kotlin, CPython para Python ---, gerando uma representação intermediária
(programa GOTO). A execução simbólica converte o GOTO em fórmulas que representam
as execuções até a profundidade $k$, das quais se geram as condições de
verificação (VCCs). As VCCs, após simplificação, vão ao solver SMT (aqui, o
Boolector~3.2.4): SAT indica um contraexemplo; UNSAT, que a propriedade vale até
$k$~\cite{Menezes2024}.

\begin{figure}[ht]
  \centering
  \resizebox{\textwidth}{!}{%
  \begin{tikzpicture}[
    font=\linespread{1}\footnotesize, >={Stealth[length=2mm]},
    bloco/.style={draw=palEdge, fill=palFill, rounded corners=2pt, align=center,
      inner sep=4pt, minimum height=9mm},
    fe/.style={draw=palEdge, fill=palFill, rounded corners=1.5pt, align=center,
      text width=2.3cm, inner sep=2pt, minimum height=8mm},
    sat/.style={draw=palEdgeDark, line width=0.8pt, fill=palStrong, rounded corners=2pt, align=center, text width=2.6cm, inner sep=3pt},
    unsat/.style={draw=palEdge, fill=palPaper, rounded corners=2pt, align=center, text width=2.6cm, inner sep=3pt}]
    \node[fe] (c) {\textbf{C/C++}\\Clang (AST)};
    \node[fe, below=1.5mm of c] (k) {\textbf{Java/Kotlin}\\Soot (Jimple)};
    \node[fe, below=1.5mm of k] (s) {\textbf{Solidity}\\gramática};
    \node[fe, below=1.5mm of s] (p) {\textbf{Python}\\AST do CPython};
    \node[above=1mm of c, font=\small\bfseries] (fet) {\textit{Front-ends}};
    \begin{scope}[on background layer]
      \node[draw=palEdge, fill=palPaper, rounded corners=3pt, fit=(fet)(c)(p), inner sep=3pt] (fes) {};
    \end{scope}
    \node[bloco, right=9mm of fes, text width=2.5cm] (goto) {\textbf{Programa GOTO}\\representação intermediária};
    \node[bloco, right=7mm of goto, text width=2.7cm] (sym) {\textbf{Execução simbólica}\\BMC, $k$-indução,\\interpretação abstrata,\\\textit{slicing}};
    \node[bloco, right=7mm of sym, text width=2.6cm] (vcc) {\textbf{Condições de verificação}\\(VCCs)};
    \node[bloco, above=8mm of vcc, text width=4.2cm, fill=palStrong, draw=palEdgeDark] (prop) {\textbf{Propriedades}\\limites de arranjos, ponteiros,\\\textit{overflow}, divisão por zero,\\\textit{assertions} do usuário};
    \node[bloco, right=7mm of vcc, text width=2.2cm] (smt) {\textbf{Solver SMT}\\Boolector, Z3,\\Bitwuzla, cvc5,\\MathSAT, Yices};
    \node[sat, right=9mm of smt, yshift=9mm] (sa) {\textbf{SAT}\\contraexemplo\\(violação no modelo)};
    \node[unsat, right=9mm of smt, yshift=-9mm] (un) {\textbf{UNSAT}\\propriedade válida\\até o limite $k$};
    \draw[->] (fes.east |- goto) -- (goto);
    \draw[->] (goto) -- (sym);
    \draw[->] (sym) -- (vcc);
    \draw[->] (prop) -- (vcc);
    \draw[->] (vcc) -- (smt);
    \draw[->] (smt.east) -- ++(4mm,0) |- (sa.west);
    \draw[->] (smt.east) -- ++(4mm,0) |- (un.west);
  \end{tikzpicture}}
  \caption{Pipeline de verificação do ESBMC: o código-fonte é processado pelos
           \textit{front-ends} correspondentes, transformado em programa GOTO e submetido
           à execução simbólica para geração das condições de verificação, que
           são resolvidas por um solver SMT: SAT indica um contraexemplo
           (violação no modelo); UNSAT, que a propriedade vale até o limite $k$.
           Elaborado pela autora com base em~\cite{ESBMC_GitHub}.}
  \label{fig:esbmc_pipeline}
\end{figure}

\subsection{Classes de Propriedades Verificadas}

O ESBMC é capaz de detectar \textit{buffer overflows}, desreferência de ponteiros
nulos, \textit{use-after-free} e vazamentos de memória; \textit{overflow} e
\textit{underflow} de inteiros; operações de ponto flutuante que produzem NaN ou
infinito; e, em programas concorrentes, \textit{data races} e
\textit{deadlocks}~\cite{Menezes2024}. Cada classe é ativada por opções próprias
da ferramenta; nesta pesquisa foram habilitadas as verificações de segurança de
memória (limites de \textit{arrays} e validade de ponteiros) e de \textit{overflow}
aritmético, conforme a configuração detalhada no Capítulo~\ref{cap:experimentos}.
O usuário pode ainda inserir \textit{assertions} próprias no código.

%% LC (resolvido): lacuna de fundamentacao -- acrescentada subsecao sobre a CWE
%% (relacao com CVE, hierarquia, e as entradas usadas no Cap. de resultados).
\subsection{Classificação de Vulnerabilidades: a CWE}

A \textit{Common Weakness Enumeration} (CWE) é uma taxonomia mantida pelo MITRE
que cataloga tipos de fraquezas de software --- padrões recorrentes de erro que
podem originar vulnerabilidades --- independentemente de qualquer ocorrência
específica~\cite{MITRE_CWE}. Ela se distingue da \textit{Common Vulnerabilities
and Exposures} (CVE), que identifica instâncias públicas e concretas em produtos
específicos e costuma referenciar uma ou mais entradas CWE para indicar a causa
raiz da falha. A CWE organiza suas entradas em uma hierarquia de generalidade, na
qual uma fraqueza mais específica é filha de uma mais abstrata; por exemplo,
CWE-787 (\textit{Out-of-bounds Write}) é filha de CWE-119 (\textit{Improper
Restriction of Operations within the Bounds of a Memory Buffer})~\cite{MITRE_CWE}.

Esta pesquisa emprega seis entradas da CWE para classificar os achados do
Capítulo~\ref{cap:resultados}: CWE-20 (\textit{Improper Input Validation}),
CWE-1284 (\textit{Improper Validation of Specified Quantity in Input}), subclasse
de CWE-20, CWE-119 (\textit{Improper Restriction of Operations within the Bounds of
a Memory Buffer}), CWE-190 (\textit{Integer Overflow or Wraparound}), CWE-191
(\textit{Integer Underflow (Wrap or Wraparound)}) e CWE-787 (\textit{Out-of-bounds
Write})~\cite{MITRE_CWE}. A nomenclatura da CWE-191 usa ``\textit{underflow}'' num
sentido específico de segurança --- uma subtração sem sinal que dá a volta para um
valor grande ---, distinto do uso em análise numérica (resultado pequeno demais
para ser representado); essa distinção é retomada no
Capítulo~\ref{cap:resultados} ao classificar o achado no GPS uBlox.

\subsection{Extensões Recentes}

Trabalhos recentes ampliaram o ESBMC ou o empregaram em novos domínios. Entre as
extensões da própria ferramenta, Monteiro et al. ampliaram o suporte a C++, com
\textit{templates} e herança virtual~\cite{Monteiro2022} --- pertinente ao PX4,
escrito em C++17. Entre os desenvolvimentos e aplicações relacionados, o FuSeBMC
combina verificação formal e geração de testes por \textit{fuzzing}, usando o ESBMC
como um de seus motores~\cite{Alshmrany2022}, e Wu et al. aplicaram o ESBMC a
componentes do Arm Confidential Computing~\cite{Wu2024}.

%% LC (pendente -- resolver ao editar o Cap. 3): o paragrafo de Extensoes Recentes
%% se repete no Cap. 3; manter uma ocorrencia e usar referencia cruzada.

\section{PX4 Autopilot}

O PX4 Autopilot é uma plataforma de software \textit{open-source} para controle de
veículos aéreos não tripulados, sob licença BSD~\cite{PX4_GitHub,
PX4_Documentation}. Suporta diversas configurações de aeronave --- multirotores,
asas fixas, VTOL e plataformas híbridas --- e é amplamente utilizado em aplicações
comerciais, científicas e recreativas.

\subsection{Arquitetura de Software}

A arquitetura do PX4 segue o padrão \textit{publish-subscribe}: os módulos
comunicam-se pelo \textit{middleware} uORB (\textit{micro Object Request Broker}),
que desacopla os componentes e facilita o desenvolvimento modular e a
manutenção~\cite{PX4_Documentation}. Os principais subsistemas incluem os
\textit{drivers} de sensores (GPS, IMU, magnetômetro e barômetro), a estimação de
estado por filtros de Kalman estendidos, os controladores de atitude e posição, a
navegação com planejamento de trajetória e os mecanismos de \textit{failsafe},
além da interface de comunicação MAVLink~\cite{PX4_Documentation}. A
Figura~\ref{fig:px4_arquitetura} resume essa organização.

\begin{figure}[ht]
  \centering
  \resizebox{0.95\textwidth}{!}{%
  \begin{tikzpicture}[
    font=\footnotesize, >={Stealth[length=2mm]},
    verif/.style={draw=palEdgeDark, fill=palStrong, line width=0.9pt, rounded corners=3pt,
      align=center, text width=3cm, minimum height=11mm, inner sep=3pt},
    plano/.style={draw=palEdge, fill=palPaper, rounded corners=3pt,
      align=center, text width=3cm, minimum height=11mm, inner sep=3pt},
    tabp/.style={draw=palEdge, fill=palFill, rounded corners=2pt,
      align=center, text width=2.1cm, minimum height=4.5mm, inner sep=1pt,
      font=\scriptsize\bfseries},
    tabs/.style={draw=palEdge, fill=palFill, rounded corners=2pt,
      align=center, text width=2.1cm, minimum height=4.5mm, inner sep=1pt,
      font=\scriptsize\bfseries},
    bus/.style={draw=palEdgeDark, fill=palFill, line width=1pt, rounded corners=4pt,
      align=center, text width=2.6cm, minimum height=18mm},
    grp/.style={draw=palEdge, dashed, rounded corners=5pt, inner sep=6pt}]

    \node[verif] (imu) {\textbf{IMU Driver}\\(BMI088)};
    \node[verif, below=10mm of imu] (gps) {\textbf{GPS Driver}\\(uBlox)};
    \node[verif, below=10mm of gps] (rc) {\textbf{RC Input / Expo}};
    \node[tabp, above=1.5mm of imu] (timu) {Publisher};
    \node[tabp, above=1.5mm of gps] (tgps) {Publisher};
    \node[tabp, above=1.5mm of rc]  (trc)  {Publisher};

    \node[bus, right=30mm of gps] (uorb) {\textbf{uORB}\\barramento de\\mensagens};

    \node[plano, right=30mm of uorb] (fc) {\textbf{Controle de voo}};
    \node[plano, above=10mm of fc]  (ekf) {\textbf{EKF / Estimador de estado}};
    \node[verif, below=10mm of fc]  (mav) {\textbf{MAVLink}};
    \node[tabs, above=1.5mm of ekf] (tekf) {Subscriber};
    \node[tabs, above=1.5mm of fc]  (tfc)  {Subscriber};
    \node[tabs, above=1.5mm of mav] (tmav) {Subscriber};

    \begin{scope}[on background layer]
      \node[grp, fit=(timu)(imu)(trc)(rc)]   (gpub) {};
      \node[grp, fit=(tekf)(ekf)(tmav)(mav)] (gsub) {};
    \end{scope}
    \node[font=\bfseries, above=1.5mm of gpub] {Publishers};
    \node[font=\bfseries, above=1.5mm of gsub] {Subscribers};

    \draw[->] (imu.east) -- node[above, font=\scriptsize\itshape] {orb\_publish()} (uorb.west);
    \draw[->] (gps.east) -- (uorb.west);
    \draw[->] (rc.east)  -- (uorb.west);
    \draw[->] (uorb.east) -- (ekf.west);
    \draw[->] (uorb.east) -- (fc.west);
    \draw[->] (uorb.east) -- (mav.west);

    \node[verif, text width=5.2cm, minimum height=7mm, below=26mm of uorb] (leg)
      {\textbf{Exemplos de módulos verificados}};
  \end{tikzpicture}}
  \caption{Arquitetura de comunicação do PX4 Autopilot segundo o padrão
    \textit{publish-subscribe}: \textit{drivers} de sensores (IMU BMI088, GPS uBlox)
    e módulos de entrada e telemetria (Expo RC, MAVLink) publicam e assinam tópicos
    por meio do \textit{middleware} uORB. Com realce (tom mais escuro e borda mais
    espessa), exemplos de módulos verificados neste trabalho; a figura é
    ilustrativa e não representa a
    totalidade dos módulos analisados (a amostra completa consta do
    Capítulo~\ref{cap:metodologia}). Elaborado pela autora com base
    em~\cite{PX4_Documentation}.}
  \label{fig:px4_arquitetura}
\end{figure}

\section{Ambiente de Execução}
\label{sec:ambiente-execucao}

O PX4 roda sobre o NuttX, sistema operacional de tempo real compatível com POSIX
e voltado a microcontroladores de recursos limitados~\cite{NuttX}, com
escalonamento preemptivo e primitivas de sincronização. Os alvos considerados
neste trabalho, baseados em microcontroladores STM32, dispõem de uma unidade de
proteção de memória (MPU), mas essa proteção só é ativada nos modos
\textit{protected} ou \textit{kernel} do NuttX; como foram compilados no modo
\textit{flat}, não há isolamento entre os módulos --- todos compartilham o mesmo
espaço de endereçamento, de modo que um defeito de memória em qualquer módulo pode
comprometer o estado global~\cite{NuttX}.
%% LC (resolvido): STM32 possui MPU -- esclarecido que ela so e usada nos modos
%% protected/kernel; nos alvos deste trabalho (modo flat) nao ha isolamento.

\subsection{Desafios para Verificação Formal}

O PX4 reúne características que tornam sua verificação ao mesmo tempo desafiadora e
necessária. O uso extensivo de ponto flutuante em navegação e controle exige
modelagem precisa da semântica IEEE-754, sob risco de não detectar defeitos
ligados a NaN ou a imprecisões numéricas; os laços de controle executam a taxas da
ordem de centenas de hertz, impondo restrições temporais rigorosas; e a interação
com hardware sujeito a falhas --- sensores que retornam dados corrompidos,
\textit{timeouts} de comunicação --- exige robustez adicional do
software~\cite{Wanner2024}.

Defeitos em sistemas de controle de voo podem resultar na perda do equipamento
ou, em cenários graves, em danos a terceiros~\cite{Siddiqi2022, Mekdad2023}. A
certificação de software aeronáutico apoia-se tradicionalmente em atividades de
teste (DO-178C); o suplemento de métodos formais DO-333 reconhece a verificação
formal como meio de conformidade: conforme os objetivos e critérios aplicáveis, a
análise formal pode cumprir parte dos objetivos de verificação e, em determinados
casos, substituir revisões, análises ou testes específicos~\cite{RTCA_DO178C,
RTCA_DO333}. Trabalhos
anteriores em sistemas embarcados \textit{safety-critical} estabeleceram
metodologias aplicáveis a VANTs~\cite{Mzid2024}, e a combinação de Model Checking
com análise de código embarcado tem se mostrado eficaz na detecção de defeitos que
escapam a abordagens baseadas apenas em testes~\cite{Khairullah2022, Sousa2026}.
%% LC (resolvido): alegacao de certificacao reformulada com o suplemento de
%% metodos formais DO-333 (nao substitui testes automaticamente). LC (resolvido):
%% Chatterjee2022 nao sustentava a frase -- trocado por Sousa2026.

\section{Resumo}

Este capítulo reuniu os fundamentos do trabalho. A verificação formal foi
apresentada como complemento aos testes, com foco em propriedades de
\textit{safety}, cuja violação, no domínio de voo, pode ser catastrófica. Entre as
técnicas, adotou-se o Model Checking, pela automação e pela geração de
contraexemplos; o \textit{Bounded Model Checking}, por limitar a busca a $k$
passos, é \textit{sound} e completo apenas até esse limite --- restrição atenuada
pela $k$-indução, aplicada adiante ao GPS uBlox. A expressividade dos solvers SMT
(\textit{bit-vectors}, \textit{arrays} e IEEE-754) é o que viabiliza verificar
firmware embarcado com precisão, e o ESBMC consolida essas capacidades. Por fim, a
arquitetura modular do PX4, articulada pelo \textit{middleware} uORB, é o que
torna viável a verificação isolada de cada módulo --- premissa central da
metodologia desenvolvida nos próximos capítulos.
